
\chapter{持续医学图像分割的场景定义与综合评测研究}
\label{chap:benchmark}
\setlength{\emergencystretch}{1em}

\section{引言}
\label{sec:benchmark-introduction}

医学图像分割是影像定量分析、治疗规划和疗效评估中的基础环节。过去十余年，深度学习显著提高了多类器官和病灶的自动分割性能\cite{litjens2017survey}，但大多数模型仍遵循一次性离线训练范式：训练集在建模前被集中整理，模型部署后参数保持固定，后续数据变化通常通过重新收集历史数据并整体训练来处理。真实临床环境并不静止。新的医疗中心、扫描设备和采集方案会改变输入分布，新的解剖结构或病灶类别会扩展输出空间，不同器官和模态也可能按任务序列逐步进入分析流程。若模型不能在这些变化中持续更新，其有效性会随数据环境演化而下降；若直接使用新任务数据进行微调，又可能覆盖已有表征并造成灾难性遗忘\cite{lee2020clinicalcl,vokinger2021medicaldevices}。

持续学习为上述问题提供了一种顺序建模思路：模型在任务或数据分布连续到达时更新参数，同时尽量保持先前获得的能力\cite{wang2024continualsurvey}。这一需求在医学影像中尤为突出。一方面，受患者隐私、数据治理和跨中心协作条件限制，历史原始数据未必能够长期完整保留或重新集中；另一方面，医学分割属于像素级预测，标签空间、背景语义、输出头和解剖结构之间的关系都会影响新旧知识如何冲突。因此，持续医学图像分割不能只套用自然图像分类中的任务定义和平均准确率评价，而需要同时明确临床变化的来源、历史信息访问方式、输出组织和测试条件。

从临床数据演化方式看，至少存在三类具有不同技术含义的变化。第一，同一分割目标在不同中心、设备或采集方案下依次到达，标签语义保持不变而输入分布发生偏移；第二，同一成像区域中的解剖类别逐步扩展，统一输出空间随阶段增加，并伴随背景语义变化；第三，器官、模态和分割目标同时改变，任务之间的标签集合可能互不重叠。这些变化分别对应本章定义的 Domain-CL、Class-CL 和 Organ-CL。三者都要求顺序学习，但可比较的性能矩阵条目、输出头形式和适用指标并不相同。

现有持续学习研究通常从域增量、类增量或任务增量等角度描述顺序学习问题，但不同医学分割工作在场景命名、标签设定、输出头、历史数据访问和测试方式上并不一致\cite{wang2024continualsurvey,kumari2025medicalcl,yuan2024continualseg,gonzalez2023lifelong,whatwrong}。相同的“跨中心分割”可能被组织为顺序域学习，也可能在训练阶段共同使用多个中心；新增器官既可能被视为类别扩展，也可能采用任务特定输出头。若这些条件没有被明确区分，方法间即使报告相同指标，其结果也未必具有可比性。

评价维度的不完整会进一步影响结论。最终阶段的平均 Dice 可以反映模型结束训练后的总体分割能力，却不能区分该结果来自历史任务保持还是当前任务学习。两个最终平均 Dice 相同的模型，可能一个较好地保持了旧任务，另一个只是最后一个任务学得更充分。遗忘量能够描述旧任务性能下降，但过强的历史保护也可能压低当前任务性能。参数隔离方法可以减少遗忘，却以模型容量增长为代价；回放方法利用旧样本缓解知识损失，但引入额外存储和历史数据访问。对于尚未参与训练的未来域，模型是否形成了可直接迁移的表示，也不能由最终平均性能和遗忘单独刻画。因而，面向临床部署的评价应同时覆盖稳定性、可塑性、前向泛化、参数效率和回放负担，而不是把“降低遗忘”视为唯一目标。

基于上述问题，本章构建持续医学图像分割的场景定义与综合评测框架。框架按照输入分布、标签空间和分割对象的变化组织 Domain-CL、Class-CL 和 Organ-CL 三类场景，并在统一网络、训练预算和数据划分下比较正则化、回放与参数隔离等代表性方法。评价指标覆盖最终性能、历史知识保持、当前任务建模、前向泛化和资源条件。其中，受限建模能力（restricted modeling ability，RMA）衡量持续约束下当前任务性能相对于独立训练参考的比例；扩展前向迁移（Extended Forward Transfer，E-FWT）汇总 Domain-CL 中模型对多个尚未参与优化域的直接表现。本章还考察任务顺序、回放容量以及基础分割模型顺序适配的影响，并讨论持续学习与联邦学习、测试时适应等相关范式的联系。所有观察均限定于相应数据、网络、访问条件和任务设定。

\begin{figure}[htbp]
\centering
\includegraphics[width=\textwidth]{ch03/benchmark_overview_new.pdf}
\caption[持续医学图像分割基准的整体构成]{持续医学图像分割基准的整体构成。基准由临床驱动的 Domain-CL、Class-CL 和 Organ-CL 场景、统一的顺序训练设置以及超越遗忘的综合评价三部分组成；评价同时考虑总体性能、稳定性、可塑性、前向泛化、参数效率和回放负担。}
\label{fig:benchmark-overview}
\end{figure}

\section{相关工作与基准定位}
\label{sec:benchmark-related-work}

\subsection{持续学习与持续语义分割综述}

通用持续学习综述通常将现有方法归纳为正则化、回放和参数隔离三类\cite{wang2024continualsurvey}。这类研究建立了灾难性遗忘、后向迁移和经验回放等基本概念，但主要证据来自图像分类。分类任务只需输出图像级标签，而医学图像分割需要保持像素级解剖边界，并处理标签空间扩展、背景语义变化和不同任务输出头等问题。因此，只使用最终准确率和后向迁移难以覆盖持续分割中的主要能力取舍。

持续语义分割综述进一步讨论了稳定性--可塑性冲突、语义漂移和背景变化等问题\cite{yuan2024continualseg}，但多数实验建立在 PASCAL VOC、ADE20K 等自然图像数据上。医学影像中的跨中心变化可能在标签空间不变时显著改变图像外观；类别增量分割受到解剖结构关系和空间先验影响；跨器官任务又会同时改变模态、结构和输出语义。这些条件使自然图像持续分割的结论不能直接等同于医学场景。

面向医学影像的持续学习综述覆盖分类、检测和分割等任务\cite{kumari2025medicalcl}，为该领域提供了更完整的文献脉络，但其目标是总结方法进展，并未针对持续医学分割建立统一的临床场景分类、数据组织和超越遗忘的实验设置。本章因此将重点放在可复现的场景定义和统一实证比较，而不是再次给出一般性方法综述。

\subsection{持续医学图像分割方法与现有基准}

已有持续医学分割方法首先集中在跨扫描仪或跨中心适应。Karani 等通过域特定归一化参数支持脑 MR 分割模型在不同扫描方案间顺序更新\cite{karani2018lifelong}，S$^3$R 则利用形状和语义选择性正则化缓解多中心持续分割中的遗忘\cite{zhang2023s3r}。另一类工作研究解剖类别或器官任务的逐步扩展，例如腹部多器官与肿瘤持续分割\cite{zhang2023continualorgan}。这些工作针对特定场景提出方法，能够说明局部问题的可行性，但由于数据划分、网络、历史信息访问和评价指标不同，跨论文比较仍较困难。

在标准化评测方面，Lifelong nnU-Net 提供了基于 nnU-Net 的顺序训练与测试框架，是持续医学分割可复现研究的重要基础\cite{gonzalez2023lifelong}，但其主要聚焦跨中心 Domain-CL，评价以 Dice 和后向迁移为主。Gonzalez 等进一步质疑复杂持续学习方法相对于多模型方案的实际收益，并提出由独立分割器与自动编码器构成的替代方案\cite{whatwrong}；该研究同样主要围绕域变化展开。现有工作由此尚未同时覆盖 Domain-CL、Class-CL 和 Organ-CL，也较少在同一设置下联合评价可塑性、前向泛化、参数增长和回放负担。

表\ref{tab:benchmark-related-work}概括了本章与相关综述和基准的差异。表中的“实验”表示该项具有实验或明确指标支持，“讨论”表示仅在概念层面涉及，“--”表示未覆盖。本章的定位不是替代上述综述或方法，而是在统一任务定义、数据划分、主干网络和评价口径下，观察不同策略在三类临床变化中的共同取舍。

\begin{table}[htbp]
\centering
\caption{持续医学图像分割相关综述与评测基准的覆盖范围比较}
\label{tab:benchmark-related-work}
\scriptsize
\renewcommand{\arraystretch}{1.12}
\setlength{\tabcolsep}{2.8pt}
\resizebox{\textwidth}{!}{%
\begin{tabular}{lcccccccccc}
\toprule
\multirow{2}{*}{工作} & \multirow{2}{*}{类型} & \multicolumn{3}{c}{场景} & \multicolumn{5}{c}{评价维度} & \multirow{2}{*}{统一设置} \\
\cmidrule(lr){3-5}\cmidrule(lr){6-10}
 & & D-CL & C-CL & O-CL & 稳定性 & 可塑性 & 前向泛化 & 参数 & 回放 & \\
\midrule
Yuan 等\cite{yuan2024continualseg} & 自然图像综述 & 讨论 & 讨论 & -- & -- & -- & -- & -- & -- & -- \\
Kumari 等\cite{kumari2025medicalcl} & 医学影像综述 & 讨论 & 讨论 & 讨论 & -- & -- & -- & -- & -- & -- \\
Lifelong nnU-Net\cite{gonzalez2023lifelong} & 评测框架 & 实验 & -- & -- & 实验 & -- & -- & -- & -- & 实验 \\
Gonzalez 等\cite{whatwrong} & 基准研究 & 实验 & -- & -- & 实验 & -- & -- & -- & -- & 实验 \\
本章 & 综合基准 & 实验 & 实验 & 实验 & 实验 & 实验 & 实验 & 实验 & 实验 & 实验 \\
\bottomrule
\end{tabular}%
}
\end{table}

\section{持续医学图像分割的问题定义}
\label{sec:benchmark-problem-definition}

设持续学习过程包含按顺序到达的任务集合
$\mathcal{T}=\{1,\ldots,T\}$。第 $t$ 个阶段获得训练集
$\mathcal{D}_t=\{(x_i^t,y_i^t)\}_{i=1}^{N_t}$，其中 $x_i^t$ 为医学图像，$y_i^t$ 为对应分割标签，$N_t$ 为该阶段的训练样本数。令 $\mathbb{D}_t$ 表示定义在图像--标签空间 $\mathcal{X}_t\times\mathcal{Y}_t$ 上的联合概率分布，则
\begin{equation}
\mathcal{D}_t=\{(x_i^t,y_i^t)\}_{i=1}^{N_t},
\qquad
(x_i^t,y_i^t)\overset{\mathrm{i.i.d.}}{\sim}\mathbb{D}_t.
\label{eq:benchmark-task-distribution}
\end{equation}
令 $Y_t\subseteq\mathcal{Y}_t$ 表示第 $t$ 阶段实际使用的标签集合。不同场景由 $\mathbb{D}_t$、$Y_t$ 以及输出组织方式随阶段的变化共同决定。这里的任务划分不是对静态数据集的任意切片，而应对应中心、类别、器官或模态等具有实际含义的变化。

完成第 $t$ 个任务训练后，模型参数记为 $\theta_t$。令 $d_{t,i}$ 表示模型 $\theta_t$ 在第 $i$ 个任务测试集上的平均 Dice，则全部可评测结果可以组织为阶段—任务性能矩阵
$\mathbf{D}=[d_{t,i}]_{t,i=1}^{T}$。矩阵的第 $t$ 行描述模型完成第 $t$ 阶段训练后的能力，第 $i$ 列描述任务 $i$ 的性能随后续更新的变化。历史任务条目用于分析知识保持，主对角线 $d_{t,t}$ 用于分析当前任务学习；在输出语义保持一致的 Domain-CL 中，位于主对角线上方的未来域条目还可用于分析前向泛化。

持续学习方法对历史信息的使用方式并不相同。正则化方法不直接保存旧图像，而是约束参数或旧模型输出的变化；回放方法保存旧样本或其他历史表示，在后续阶段重新参与训练；参数隔离方法为不同任务保留独立或部分独立的模型容量。为了区分算法能力与访问条件，本章还设置两类参照：Non-CL 按任务顺序直接更新模型，不使用专门的抗遗忘机制；JointTrain 将全部任务数据集中训练。JointTrain 能够访问完整历史数据，因此只作为不同访问条件下的性能参照，不作为与受限持续学习方法完全同条件的比较对象。

测试方式还需结合输出空间解释。Domain-CL 保持标签语义不变，使用共享输出头；Class-CL 逐步扩展统一标签空间，并要求最终模型预测全部已见类别；Organ-CL 的任务目标、成像模态和输出语义均可改变，因而采用任务特定输出头。本章采用任务边界已知的阶段式训练，并按上述输出头完成测试。该设定不涉及任务身份未知（task-free）的路由、任务边界检测或开放世界识别。


\section{持续医学图像分割基准场景}
\label{sec:benchmark-scenarios}

三类场景共享顺序训练和历史能力保持这一基本要求，但输入分布、标签空间和输出组织方式不同。Domain-CL 关注相同分割语义在不同中心或设备域之间的持续适应；Class-CL 关注同一图像域内分割类别的逐步扩展；Organ-CL 则面向器官、病灶和模态均可能改变的任务序列。表\ref{tab:benchmark-scenarios}概括了三类场景的主要区别。

\begin{table}[htbp]
\centering
\caption{持续医学图像分割三类基准场景的定义与实验设定差异}
\label{tab:benchmark-scenarios}
\small
\renewcommand{\arraystretch}{1.18}
\resizebox{\textwidth}{!}{%
\begin{tabular}{p{2.0cm}p{2.8cm}p{2.8cm}p{3.0cm}p{3.2cm}}
\toprule
场景 & 输入分布与标签空间 & 输出组织 & 任务序列实例 & 主要评价重点 \\
\midrule
Domain-CL &
$P(\mathcal{X}_{t_1})\neq P(\mathcal{X}_{t_2})$，$Y_{t_1}=Y_{t_2}$ &
共享二通道输出头，分割语义保持不变 &
六个前列腺 MRI 中心按顺序到达 &
跨域知识保持、当前域学习、E-FWT 与资源开销 \\
\addlinespace
Class-CL &
图像域相对稳定，$Y_{t_1}\neq Y_{t_2}$，统一标签空间逐步扩展 &
共享八通道输出头，最终预测全部已见结构 &
七个心脏结构分三阶段加入 &
背景语义变化、旧类保持、新类学习、WCD 与资源开销 \\
\addlinespace
Organ-CL &
$P(\mathcal{X}_{t_1})\neq P(\mathcal{X}_{t_2})$，且 $Y^{+}_{t_1}\cap Y^{+}_{t_2}=\varnothing$ &
每个任务使用二通道任务特定输出头 &
左心房 MRI、前列腺 MRI、肝 CT、脑肿瘤 FLAIR MRI &
跨任务知识保持、当前任务学习与参数扩展 \\
\bottomrule
\end{tabular}
}
\end{table}

\subsection{Domain-CL：跨中心与设备域增量分割}
\label{subsec:benchmark-domain-cl}

Domain-CL 描述分割目标保持不变，而图像分布随中心、设备或采集参数发生变化的情况。对任意两个任务 $t_1$ 和 $t_2$，其标签集合满足 $Y_{t_1}=Y_{t_2}$，但输入分布通常满足
$P(\mathcal{X}_{t_1})\neq P(\mathcal{X}_{t_2})$。本章使用六个中心的前列腺 T2 加权 MRI 构建任务序列，所有阶段均预测前列腺与背景，因此网络使用共享的二通道输出头。

该场景与静态多域训练的区别在于，不同中心的数据并非一次性进入训练过程。模型在学习后续中心时仍需保持此前中心的分割能力。它也不同于只在目标域适配的域适应问题，因为评价不仅关注当前中心，还跟踪完整历史域的性能。由于不同域具有相同输出语义，模型在尚未训练域上的直接 Dice 具有可比较性，因此 Domain-CL 还可以评价前向泛化。

跨中心差异可能来自扫描仪、采集参数、人群构成和数据预处理。基准并不假设这些变化能够被单一域标签完全解释，而是把中心序列作为可复现的域增量设定。三类场景也不意味着所有医学域变化都可归入该设定；连续漂移、中心重访或多个变化因素叠加发生的情况仍超出当前实验范围。

\subsection{Class-CL：类别增量分割}
\label{subsec:benchmark-class-cl}

Class-CL 保持成像区域和图像分布相对稳定，但分割类别随阶段扩展。设第 $t$ 阶段提供标签集合 $Y_t$，则不同阶段满足 $Y_{t_1}\neq Y_{t_2}$，最终模型需要在统一输出空间中预测全部已见类别。本章使用心脏 CT 数据，将七个解剖结构依次组织为三个阶段：第一阶段学习左心室、左心房和心肌，第二阶段加入右心室和右心房，第三阶段加入升主动脉和肺动脉。

类别增量分割不仅涉及输出维度扩展，还受到背景语义变化影响。当前阶段只对新增类别提供监督时，旧类别或未来类别像素可能被并入背景，导致“背景”在不同阶段不再表示相同语义。普通知识蒸馏或参数约束若忽略这一变化，可能把当前新增结构误当作旧背景，或使旧类别在后续阶段被覆盖。基准因此引入 MiB 和 PLOP 等面向持续语义分割的专项方法作为对照\cite{cermelli2020mib,douillard2021plop}。

Class-CL 采用共享八通道输出头，最终阶段使用全类别 Dice 评价统一分割能力。训练与评价遵循上述三阶段标签设定，并据此分析背景语义变化和旧类别保持。相关结论适用于该输出头与标签可访问设置，不直接代表其他类别增量标注方式。

\subsection{Organ-CL：跨器官任务持续学习}
\label{subsec:benchmark-organ-cl}

Organ-CL 进一步改变分割对象、图像分布和成像模态。令 $Y_t^{+}=Y_t\setminus\{0\}$ 表示不含背景的前景类别集合。不同任务的前景类别互不重叠，即
$Y^{+}_{t_1}\cap Y^{+}_{t_2}=\varnothing$，输入分布也通常满足
$P(\mathcal{X}_{t_1})\neq P(\mathcal{X}_{t_2})$。本章构建四阶段任务序列：左心房 LGE MRI、前列腺 MRI、肝脏 CT 和脑肿瘤 FLAIR MRI。每个任务均使用独立二通道输出头，主干网络在任务之间共享。

该场景与 Class-CL 的差别不只是标签集合是否重叠。不同器官和模态对应不同图像外观、空间结构及分割目标，直接在尚未训练任务上计算跨任务 Dice 缺少统一语义基础。因此，E-FWT 不用于 Organ-CL。评价主要关注模型完成全部顺序任务后的平均性能、旧任务保持、当前任务建模能力和参数增长。

Organ-CL 也不同于静态多任务学习。静态多任务训练可以同时使用全部任务数据并共同优化共享表示，而本章按顺序提供任务，模型在学习新器官时不能默认完整访问旧任务训练集。多头结构减少了输出类别之间的直接竞争，但共享主干仍可能发生表征覆盖；为每个任务扩展网络还会带来额外参数开销。


\section{数据组织与统一实验设定}
\label{sec:benchmark-protocol}

\subsection*{数据集与任务划分}

三类场景中的每个分割任务均在病例层面随机划分为 60\% 训练集、15\% 验证集和 25\% 测试集，验证集用于方法超参数选择。不同持续学习方法使用相同的任务定义和数据划分，避免因训练样本或验证样本不一致影响比较。

Domain-CL 使用 NCI-ISBI 2013、I2CVB 和 PROMISE12 前列腺 MRI 数据构造六个中心域，其中 NCI-ISBI 2013 对应中心 A 和 B，I2CVB 对应中心 C，PROMISE12 对应中心 D、E 和 F\cite{dataset_prostate2,dataset_prostate,dataset_prostate3,dataset_domaincl}。中心 C 的图像先进行裁剪和空间对齐，所有中心的轴位切片统一为 $256\times256$，并依据非零体素计算均值和标准差完成归一化。六个中心按照固定顺序形成持续任务。

Class-CL 使用 MMWHS 心脏 CT 数据\cite{care25ASA1,discussion_inter_seg3,care25ASA2}。图像被重采样至 $0.78\times0.78$ mm 的平面分辨率，平均层厚为 1.60 mm，并调整为 $256\times256$。七个心脏结构分为三个阶段：左心室、左心房和心肌首先进入训练，随后加入右心室和右心房，最后加入升主动脉和肺动脉。

Organ-CL 依次使用 LAScarQS 左心房 LGE MRI、PROMISE12 中心 D 的前列腺 MRI、LiTS 肝脏 CT 和 FeTS 2021 脑肿瘤 FLAIR MRI\cite{dataset_lascars,dataset_lascars2,dataset_prostate3,dataset_taskcl,dataset_taskcl2}。各图像采用零均值和单位方差归一化，并将切片统一为 $256\times256$。

\begin{table}[htbp]
\centering
\caption{三类场景的数据组织与任务序列}
\label{tab:benchmark-datasets}
\small
\renewcommand{\arraystretch}{1.16}
\resizebox{\textwidth}{!}{%
\begin{tabular}{p{1.8cm}p{3.2cm}p{2.1cm}p{4.0cm}p{3.1cm}}
\toprule
场景 & 数据与模态 & 阶段数 & 任务顺序 & 输出与预处理 \\
\midrule
Domain-CL &
NCI-ISBI 2013、I2CVB、PROMISE12 前列腺 T2 MRI &
6 &
中心 A$\rightarrow$B$\rightarrow$C$\rightarrow$D$\rightarrow$E$\rightarrow$F &
共享二通道头；轴位 $256\times256$；非零体素归一化 \\
\addlinespace
Class-CL &
MMWHS 心脏 CT &
3 &
LV/LA/MYO $\rightarrow$ RV/RA $\rightarrow$ AA/PA &
共享八通道头；$0.78\times0.78$ mm 平面分辨率；$256\times256$ \\
\addlinespace
Organ-CL &
LAScarQS、PROMISE12 中心 D、LiTS、FeTS 2021 &
4 &
左心房 MRI $\rightarrow$ 前列腺 MRI $\rightarrow$ 肝 CT $\rightarrow$ 脑肿瘤 FLAIR MRI &
每任务二通道头；零均值单位方差归一化；$256\times256$ \\
\bottomrule
\end{tabular}
}
\end{table}

\subsection*{网络与训练设置}

为减少基础网络差异对比较的影响，所有方法采用 ResUNet32 主干。Domain-CL 使用二通道输出头，Class-CL 使用统一八通道输出头，Organ-CL 为每个任务配置二通道输出头。每个任务训练 150 个 epoch，优化器为随机梯度下降，初始学习率为 0.008，在第 80 个 epoch 后乘以 0.5，批大小为 8。Domain-CL 和 Organ-CL 使用交叉熵损失；Class-CL 使用 MiB 中用于缓解背景语义变化的无偏交叉熵。

回放方法的默认缓存容量为 32，回放批大小为 8，并采用水库采样更新缓存；Class-CL 中的回放方法按阶段标签空间进行适配。各对比方法的方法特定配置汇总于表~\ref{tab:benchmark-baseline-config}。

\begin{table}[htbp]
\centering
\caption{基准对比方法的实现配置与方法特定超参数}
\label{tab:benchmark-baseline-config}
\scriptsize
\renewcommand{\arraystretch}{1.08}
\setlength{\tabcolsep}{3.0pt}
\begin{tabular}{p{2.45cm}p{8.70cm}p{2.20cm}}
\toprule
方法 & 实现配置与方法特定超参数 & 适用范围 \\
\midrule
Non-CL、JointTrain & 无额外持续学习超参数，沿用统一的 ResUNet32、优化器、学习率、训练轮数和批大小；JointTrain 可联合访问全部任务数据。 & 三类场景 \\
Regu-EWC & 正则权重 $\lambda=1$，历史重要性衰减系数 $\gamma=0.1$。 & 三类场景 \\
Regu-SI & 路径积分代理损失权重 $c=5$，数值稳定项 $\xi=1.0$。 & 三类场景 \\
Regu-LwF & 蒸馏惩罚权重为 3，softmax 温度为 2，权重衰减为 $0.0005$。 & 三类场景 \\
Repl-ER & 缓存容量为 32，回放批大小为 8，采用水库采样。 & 三类场景 \\
Repl-GSS & 用于梯度比较的批大小为 8；其余回放设置沿用默认配置。 & 三类场景 \\
Repl-GEM & 梯度约束的 margin 参数 $\gamma=0.5$；其余回放设置沿用默认配置。 & 三类场景 \\
Repl-GPM & 构建梯度子空间的能量阈值为 $0.97$。 & 三类场景 \\
Repl-AGEM & 缓存容量为 32，回放批大小为 8，采用水库采样。 & 三类场景 \\
Repl-DER & 历史输出蒸馏损失的惩罚权重为 $0.4$；其余回放设置沿用默认配置。 & 三类场景 \\
Repl-DER++ & 采用 DER++ 目标，缓存容量为 32，回放批大小为 8，并采用水库采样。 & 三类场景 \\
Repl-FDR & 函数空间正则的惩罚权重为 $0.5$；其余回放设置沿用默认配置。 & 三类场景 \\
ParIso-PNN & 不引入额外方法超参数，按任务扩展网络分支。 & 三类场景 \\
ParIso-DAN & 采用线性组合完成任务适配。 & 三类场景 \\
ClassCL-PLOP & 池化输出蒸馏权重为 $0.01$，类别不确定性阈值为 $0.001$。 & Class-CL \\
ClassCL-MiB & 知识蒸馏损失权重为 10，无偏知识蒸馏中的真值标签权重为 $\alpha=1.0$。 & Class-CL \\
\bottomrule
\end{tabular}
\end{table}

结果表沿用现有实验记录报告均值和标准差。由于各设置未统一进行显著性检验，后文只比较数值趋势，不把方法间差异表述为统计显著。Class-CL 回放方法需要按照阶段标签空间处理缓冲区监督，其结果解释限定于本章采用的数据划分、标签设定和实现配置。

\subsection*{对比方法}

本章选择正则化、回放和参数隔离三类具有代表性的通用持续学习策略，并在 Class-CL 中增加面向持续语义分割的专项方法。选择这些方法的目的不是穷尽所有持续学习算法，而是在统一网络和训练预算下覆盖不同的历史知识载体：参数重要性、旧模型输出、原始样本、梯度约束、低维子空间以及任务特定参数。

\textbf{正则化策略。} Regu-EWC 在完成一个任务后利用 Fisher 信息近似参数重要性，并在后续任务中对重要参数偏移施加二次惩罚；Regu-SI 沿训练轨迹累计参数对损失下降的贡献，以路径积分形式估计历史重要性；Regu-LwF 保存上一阶段模型，并在当前任务样本上通过知识蒸馏约束新旧输出\cite{kirkpatrick2017ewc,method_regu_si,li2018lwf}。这类方法不保存历史原始图像，也不扩展主干参数，但其效果依赖约束强度：约束过弱时旧知识仍会被覆盖，约束过强时当前任务可塑性可能下降。

\textbf{回放与历史梯度约束策略。} Repl-ER 直接将缓冲区样本与当前小批量共同训练，并通过水库采样在固定容量下近似保留数据流中的样本；Repl-GSS 使用梯度比较提高缓冲区样本的代表性和多样性；Repl-GEM 要求当前更新不能增大记忆样本上的损失，Repl-AGEM 则以记忆样本的平均梯度构造较低开销的参考约束\cite{method_replay_er,method_replay_gss,method_replay_gem,method_replay_agem}。Repl-GPM 不回放原始图像，而是从历史任务梯度中构建主子空间，并将新任务更新投影到受保护方向的正交补\cite{saha2021gpm}。Repl-DER 在缓冲区中保存历史样本的模型输出并使用蒸馏约束，DER++ 进一步联合历史标签监督；Repl-FDR 则在记忆样本上约束新旧模型的函数距离\cite{method_replay_der,method_replay_xder,method_replay_fdr}。除 GPM 等仅保留统计表示的方法外，多数回放策略需要保存历史原始图像，这在临床场景中会引入存储、治理和隐私方面的额外条件。

\textbf{参数隔离策略。} ParIso-PNN 为后续任务增加新的网络分支，并通过横向连接利用已有分支的表示；ParIso-DAN 保留基础网络，通过控制模块对新增参数与历史参数进行线性组合以适配新任务\cite{method_dyn_pnn,method_dyn_dan}。任务特定参数减少了新旧任务在同一参数上的直接竞争，因此通常能够获得较小遗忘，但模型规模会随任务数量增长。MPE 用于量化这一代价。

\textbf{类别增量专项策略与参照。} ClassCL-MiB 通过无偏交叉熵和无偏知识蒸馏重新组织背景概率，减少旧类别在后续阶段被吸收到背景中的情况；ClassCL-PLOP 使用池化输出蒸馏和伪标签恢复旧类别监督\cite{cermelli2020mib,douillard2021plop}。二者仅在 Class-CL 中使用。Non-CL 表示不采用专门抗遗忘机制的顺序更新，JointTrain 则在同一模型中联合访问全部任务数据。后者违反受限历史访问条件，只用于表示集中训练参照。

\begin{table}[htbp]
\centering
\caption{基准中比较的方法类别及其主要历史信息载体}
\label{tab:benchmark-methods}
\small
\renewcommand{\arraystretch}{1.15}
\resizebox{\textwidth}{!}{%
\begin{tabular}{p{2.0cm}p{4.0cm}p{4.0cm}p{3.6cm}}
\toprule
类别 & 代表方法 & 历史知识载体或约束方式 & 主要附加代价 \\
\midrule
正则化 &
Regu-EWC、Regu-SI、Regu-LwF &
参数重要性、路径积分或旧模型输出 &
不保存原始样本，固定参数规模；保护过强可能限制当前任务学习 \\
\addlinespace
回放 &
Repl-ER、GSS、GEM、GPM、AGEM、DER、DER++、FDR &
原始样本、梯度约束、旧模型输出或低维子空间 &
缓存、历史统计和附加优化开销；不同方法的回放对象不同 \\
\addlinespace
参数隔离 &
ParIso-PNN、ParIso-DAN &
任务特定网络分支或适配参数 &
遗忘较小，但模型参数随任务增加 \\
\addlinespace
类别增量专项 &
ClassCL-PLOP、ClassCL-MiB &
背景修正、蒸馏和伪标签 &
依赖 Class-CL 标签设定，不能直接外推至其他场景 \\
\addlinespace
参照 &
Non-CL、JointTrain &
无专门保护；完整联合访问 &
JointTrain 使用完整历史数据，访问条件不同 \\
\bottomrule
\end{tabular}
}
\end{table}


\section{综合评价指标}
\label{sec:benchmark-metrics}

设 $d_{t,i}$ 为完成第 $t$ 个任务训练后的模型在任务 $i$ 测试集上的平均 Dice。传统评价通常只报告最终模型在所有任务上的平均性能或遗忘量。本章保留这两类指标，并进一步引入当前任务建模、前向泛化和资源开销指标。不同指标对应性能矩阵中的不同位置，不能相互替代。

\subsection*{最终平均性能}

最终平均 Dice（average Dice，A-Dice）计算模型完成全部 $T$ 个任务后在所有任务上的平均性能：
\begin{equation}
\mathrm{A\mbox{-}Dice}
=
\frac{1}{T}\sum_{i=1}^{T}d_{T,i}.
\label{eq:benchmark-adice}
\end{equation}
A-Dice 越大，表示最终模型在全部任务上的总体分割性能越高。该指标会把不同任务压缩为一个平均值，无法区分历史保持和当前任务学习，也可能受到某些高分任务的影响。

\subsection*{后向迁移与遗忘}

后向迁移率（backward transfer rate，BWTR）比较最终模型在旧任务上的性能与刚完成该任务训练时的性能：
\begin{equation}
\mathrm{BWTR}
=
\frac{1}{T-1}
\sum_{i=1}^{T-1}
\frac{d_{T,i}-d_{i,i}}{d_{i,i}}.
\label{eq:benchmark-bwtr}
\end{equation}
当 $d_{i,i}>0$ 时，BWTR 为 0 表示任务 $i$ 的最终性能与初学完成时相同；负值表示性能下降，数值越小通常对应更强遗忘；正值则表示后续任务学习改善了旧任务性能。该指标使用相对变化，因而不同任务的初始 Dice 会影响归一化结果。

\subsection*{受限建模能力}

仅限制遗忘可能使模型过度保护旧参数，而无法充分学习当前任务。为此，本章提出受限建模能力 RMA。令 $d_i^{NC}$ 表示只在任务 $i$ 上独立训练的 Non-CL 模型性能，则
\begin{equation}
\mathrm{RMA}
=
\frac{1}{T-1}
\sum_{i=2}^{T}
\frac{d_{i,i}}{d_i^{NC}}.
\label{eq:benchmark-rma}
\end{equation}
RMA 排除第一个任务，因为此时尚未受到历史任务约束。RMA 接近 1 表示持续模型在当前任务上的性能接近独立训练参考；小于 1 表示历史保护或顺序训练限制了当前任务建模；大于 1 则表示在该任务上共享历史表征带来了高于独立训练的结果。RMA 只评价当前任务学习，不能说明旧任务是否被保持。

\subsection*{扩展前向迁移与前向泛化}

本章统一使用前向泛化描述模型在尚未参与优化的未来域上的直接表现。令 $d_{0,i}$ 表示随机初始化模型在任务 $i$ 上的 Dice，则 Domain-CL 的扩展前向迁移指标定义为
\begin{equation}
\mathrm{E\mbox{-}FWT}
=
\frac{2}{T(T-1)}
\sum_{t=1}^{T-1}
\sum_{i=t+1}^{T}
\left(d_{t,i}-d_{0,i}\right).
\label{eq:benchmark-efwt}
\end{equation}
该指标汇总每个已完成阶段对所有未来域的直接表现，而不是只比较相邻任务。E-FWT 越大，说明当前模型在尚未训练域上的前向泛化相对于随机初始化越强。由于 Class-CL 的未来类别尚未进入输出语义，Organ-CL 的任务输出互不相同，E-FWT 仅用于 Domain-CL。其数值反映当前基准下的前向泛化，不应被解释为未知临床数据上的通用保证。

\subsection*{全类别分割性能}

Class-CL 最终需要在统一标签空间中预测所有已见结构。本章使用全类别 Dice（whole-class Dice，WCD）记录最终模型对全部类别的分割结果：
\begin{equation}
\mathrm{WCD}=d_{\mathrm{all}}.
\label{eq:benchmark-wcd}
\end{equation}
WCD 越大，表示最终统一输出空间中的全类别分割性能越高。它与 A-Dice 的汇总对象不同，因此两者可能给出不同方法排序。

\subsection*{模型与数据资源}

模型参数扩展率（model parameter expansion，MPE）描述任务到达后新增参数相对于初始模型参数量的平均比例：
\begin{equation}
\mathrm{MPE}
=
\frac{1}{T-1}
\sum_{t=2}^{T}
\frac{
\mathrm{Param}(\theta_t)-\mathrm{Param}(\theta_{t-1})
}{
\mathrm{Param}(\theta_1)
}.
\label{eq:benchmark-mpe}
\end{equation}
MPE 越小，表示模型容量随任务增长越少。只有初始主干和参数统计口径一致时，该指标才适合横向比较。

数据回放率（data replay ratio，DRR）衡量持续训练中回放数据相对于当前训练数据的平均比例：
\begin{equation}
\mathrm{DRR}
=
\frac{1}{T-1}
\sum_{t=1}^{T-1}
\frac{
N_t^{\mathrm{replay}}
}{
N_t^{\mathrm{train}}
}.
\label{eq:benchmark-drr}
\end{equation}
其中，$N_t^{\mathrm{replay}}$ 表示以历史任务 $t$ 为来源、在后续阶段重新使用或用于构造回放信息的数据量，$N_t^{\mathrm{train}}$ 表示任务 $t$ 的训练数据量。DRR 越小，表示对历史信息回访的依赖越低。表中带星号的 DRR 表示回放原始图像；不带星号的非零值可对应其他历史表示。MPE 和 DRR 只描述部分资源条件，不能替代计算时间、显存和通信开销的完整核算。

上述指标从最终总体性能、历史保持、当前任务学习、未来域能力和资源使用等方面观察持续模型。评价时应联合报告，而不能只选择对某一方法有利的单一指标。


\section{实验结果与分析}
\label{sec:benchmark-results}

\subsection*{Domain-CL 结果}

表\ref{tab:benchmark-domain-results}给出了六中心前列腺分割序列上的综合结果。JointTrain 的 A-Dice 为 $0.830\pm0.040$，高于所有持续学习方法，但该参照可以访问全部中心数据，不能视为相同历史访问条件下的比较结果。参数隔离方法 ParIso-PNN 和 ParIso-DAN 的 A-Dice 分别为 $0.789\pm0.002$ 和 $0.786\pm0.008$，BWTR 均为 0，说明在当前输出组织下任务特定参数能够避免最终旧域性能下降。其代价体现在 MPE：PNN 为 1.112，DAN 为 0.111。

在固定参数规模的方法中，Repl-ER 的 A-Dice 为 $0.750\pm0.021$，是表中回放类方法的较高结果；Repl-AGEM 为 $0.740\pm0.015$，其 BWTR 为 $-0.093\pm0.023$。正则化方法 EWC 和 SI 的 RMA 均高于 1，说明当前任务性能可以达到或超过独立训练参考，但其 BWTR 分别为 $-0.242\pm0.039$ 和 $-0.236\pm0.042$。这一组合表明，当前任务学习较强并不自动意味着历史知识保持良好。相反，Regu-LwF 的 BWTR 退化较小，但 RMA 仅为 $0.894\pm0.021$，说明较强的旧模型约束可能影响当前域学习。

E-FWT 进一步揭示前向泛化表现。Non-CL 的 E-FWT 为 $0.260\pm0.151$，多数正则化和回放方法位于 0.212 至 0.259 之间，没有形成稳定一致的提升。ParIso-DAN 的 E-FWT 为 $0.312\pm0.210$，是表内较高结果，但标准差也较大；ParIso-PNN 的 E-FWT 为 $0.235\pm0.148$。因此，在该中心序列中，降低遗忘与提高前向泛化不是同一目标。

\begin{table}[htbp]
\centering
\caption{Domain-CL 场景下不同持续学习方法的综合结果}
\label{tab:benchmark-domain-results}
\scriptsize
\setlength{\tabcolsep}{3.8pt}
\renewcommand{\arraystretch}{1.12}
\resizebox{\textwidth}{!}{%
\begin{tabular}{lcccccc}
\toprule
方法 & A-Dice & BWTR & RMA & E-FWT & MPE & DRR \\
\midrule
\multicolumn{7}{l}{\textit{正则化方法}} \\
Regu-EWC & $.672\pm.021$ & $-.242\pm.039$ & $1.086\pm.014$ & $.255\pm.152$ & $0$ & $0$ \\
Regu-SI & $.676\pm.024$ & $-.236\pm.042$ & $1.084\pm.008$ & $.259\pm.151$ & $0$ & $0$ \\
Regu-LwF & $.663\pm.029$ & $-.090\pm.027$ & $.894\pm.021$ & $.224\pm.153$ & $0$ & $0$ \\
\addlinespace[2pt]
\multicolumn{7}{l}{\textit{回放方法}} \\
Repl-ER & $.750\pm.021$ & $-.100\pm.027$ & $1.054\pm.014$ & $.256\pm.150$ & $0$ & $.027^{*}$ \\
Repl-GSS & $.726\pm.026$ & $-.140\pm.040$ & $1.059\pm.010$ & $.255\pm.150$ & $0$ & $.036^{*}$ \\
Repl-GEM & $.718\pm.018$ & $-.095\pm.022$ & $.994\pm.024$ & $.251\pm.154$ & $0$ & $.042^{*}$ \\
Repl-GPM & $.661\pm.010$ & $-.188\pm.010$ & $1.013\pm.022$ & $.212\pm.080$ & $0$ & $.113$ \\
Repl-AGEM & $.740\pm.015$ & $-.093\pm.023$ & $1.023\pm.027$ & $.253\pm.155$ & $0$ & $.042^{*}$ \\
Repl-DER & $.715\pm.011$ & $-.126\pm.015$ & $1.018\pm.016$ & $.249\pm.148$ & $0$ & $.027^{*}$ \\
Repl-DER++ & $.725\pm.013$ & $-.105\pm.018$ & $1.012\pm.013$ & $.251\pm.149$ & $0$ & $.027^{*}$ \\
Repl-FDR & $.728\pm.015$ & $-.111\pm.028$ & $1.025\pm.033$ & $.252\pm.136$ & $0$ & $.042^{*}$ \\
\addlinespace[2pt]
\multicolumn{7}{l}{\textit{参数隔离方法}} \\
ParIso-PNN & $.789\pm.002$ & $0\pm0$ & $1.003\pm.014$ & $.235\pm.148$ & $1.112$ & $0$ \\
ParIso-DAN & $.786\pm.008$ & $0\pm0$ & $1.016\pm.021$ & $.312\pm.210$ & $.111$ & $0$ \\
\addlinespace[2pt]
\multicolumn{7}{l}{\textit{参照方法}} \\
Non-CL & $.676\pm.021$ & $-.237\pm.038$ & $1.086\pm.008$ & $.260\pm.151$ & $0$ & $0$ \\
JointTrain & $.830\pm.040$ & -- & -- & -- & $0$ & $0$ \\
\bottomrule
\end{tabular}%
}
\vspace{2pt}
\begin{minipage}{0.98\textwidth}
\footnotesize 注：数值为均值$\pm$标准差；“--”表示不适用或未报告；DRR 中带“$^{*}$”的数值表示回放原始图像。JointTrain 可访问全部任务数据，其结果仅作为不同访问条件下的参照。
\end{minipage}
\end{table}

\subsection*{Class-CL 结果}

表\ref{tab:benchmark-class-results}展示了心脏结构逐步扩展时的结果。JointTrain 的 A-Dice 为 $0.810\pm0.023$。ParIso-PNN 的 A-Dice 为 $0.770\pm0.012$，BWTR 为 0，但 MPE 达到 1.056。ParIso-DAN 同样实现零 BWTR，其 A-Dice 为 $0.729\pm0.009$，MPE 为 0.111。参数隔离可以避免共享参数被后续类别覆盖，但不同扩展方式在性能与容量之间仍有差异。

面向背景语义变化的专项方法表现出与通用方法不同的取舍。ClassCL-MiB 的 A-Dice 为 $0.751\pm0.008$，BWTR 为 $-0.039\pm0.003$，是固定参数方法中遗忘较小的结果；ClassCL-PLOP 的 WCD 为 $0.707\pm0.009$，高于 MiB 的 $0.697\pm0.008$，但其 A-Dice 和 BWTR 分别为 $0.708\pm0.025$ 和 $-0.104\pm0.026$。WCD 与 A-Dice 的排序差异说明，最终全类别输出与逐阶段任务平均不能相互替代。

通用回放方法中，Repl-GEM 的 A-Dice 为 $0.684\pm0.036$，Repl-FDR 为 $0.666\pm0.012$。后者 BWTR 为 $-0.146\pm0.046$，优于多数回放基线，但 RMA 为 $0.917\pm0.035$。Repl-ER 的 RMA 为 $1.015\pm0.029$，说明当前任务学习接近独立训练，但 BWTR 仍为 $-0.243\pm0.068$。类别增量结果再次表明，RMA、BWTR 和 WCD 分别刻画不同能力，单一数值不能形成完整排序。

\begin{table}[htbp]
\centering
\caption{Class-CL 场景下不同持续学习方法的综合结果}
\label{tab:benchmark-class-results}
\scriptsize
\setlength{\tabcolsep}{3.8pt}
\renewcommand{\arraystretch}{1.12}
\resizebox{\textwidth}{!}{%
\begin{tabular}{lcccccc}
\toprule
方法 & A-Dice & BWTR & RMA & WCD & MPE & DRR \\
\midrule
\multicolumn{7}{l}{\textit{正则化方法}} \\
Regu-EWC & $.487\pm.020$ & $-.552\pm.042$ & $.999\pm.022$ & $.513\pm.022$ & $0$ & $0$ \\
Regu-SI & $.469\pm.007$ & $-.602\pm.016$ & $1.018\pm.028$ & $.570\pm.020$ & $0$ & $0$ \\
Regu-LwF & $.384\pm.022$ & $-.575\pm.013$ & $.782\pm.046$ & $.602\pm.002$ & $0$ & $0$ \\
\addlinespace[2pt]
\multicolumn{7}{l}{\textit{回放方法}} \\
Repl-ER & $.654\pm.034$ & $-.243\pm.068$ & $1.015\pm.029$ & $.596\pm.019$ & $0$ & $.011^{*}$ \\
Repl-GSS & $.456\pm.003$ & $-.611\pm.016$ & $.991\pm.038$ & $.600\pm.006$ & $0$ & $.018^{*}$ \\
Repl-GEM & $.684\pm.036$ & $-.304\pm.038$ & $1.014\pm.031$ & $.622\pm.031$ & $0$ & $.010^{*}$ \\
Repl-GPM & $.575\pm.023$ & $-.219\pm.033$ & $.892\pm.014$ & $.513\pm.014$ & $0$ & $.068$ \\
Repl-AGEM & $.595\pm.017$ & $-.370\pm.032$ & $1.020\pm.019$ & $.562\pm.026$ & $0$ & $.010^{*}$ \\
Repl-DER & $.628\pm.027$ & $-.219\pm.070$ & $.913\pm.029$ & $.622\pm.029$ & $0$ & $.011^{*}$ \\
Repl-DER++ & $.638\pm.036$ & $-.208\pm.056$ & $.932\pm.034$ & $.633\pm.043$ & $0$ & $.011^{*}$ \\
Repl-FDR & $.666\pm.012$ & $-.146\pm.046$ & $.917\pm.035$ & $.627\pm.015$ & $0$ & $.010^{*}$ \\
\addlinespace[2pt]
\multicolumn{7}{l}{\textit{参数隔离方法}} \\
ParIso-PNN & $.770\pm.012$ & $0\pm0$ & $.976\pm.021$ & $.620\pm.020$ & $1.056$ & $0$ \\
ParIso-DAN & $.729\pm.009$ & $0\pm0$ & $.897\pm.022$ & $.543\pm.127$ & $.111$ & $0$ \\
\addlinespace[2pt]
\multicolumn{7}{l}{\textit{类别增量专项方法}} \\
ClassCL-PLOP & $.708\pm.025$ & $-.104\pm.026$ & $.986\pm.038$ & $.707\pm.009$ & $0$ & $0$ \\
ClassCL-MiB & $.751\pm.008$ & $-.039\pm.003$ & $1.004\pm.018$ & $.697\pm.008$ & $0$ & $0$ \\
\addlinespace[2pt]
\multicolumn{7}{l}{\textit{参照方法}} \\
Non-CL & $.471\pm.012$ & $-.598\pm.015$ & $1.014\pm.028$ & $.572\pm.023$ & $0$ & $0$ \\
JointTrain & $.810\pm.023$ & -- & -- & -- & $0$ & $0$ \\
\bottomrule
\end{tabular}%
}
\vspace{2pt}
\begin{minipage}{0.98\textwidth}
\footnotesize 注：数值为均值$\pm$标准差；“--”表示不适用或未报告；DRR 中带“$^{*}$”的数值表示回放原始图像。JointTrain 可访问全部任务数据，其结果仅作为不同访问条件下的参照。
\end{minipage}
\end{table}

\subsection*{Organ-CL 结果}

表\ref{tab:benchmark-organ-results}给出了跨器官和跨模态任务序列的结果。JointTrain 的 A-Dice 为 $0.881\pm0.010$。ParIso-PNN 与 ParIso-DAN 的 A-Dice 分别为 $0.839\pm0.004$ 和 $0.824\pm0.005$，BWTR 均为 0，MPE 分别为 1.074 和 0.166。多头和任务特定参数可以隔离器官任务，但其容量增长随具体结构设计而不同。

在回放方法中，Repl-ER 的 A-Dice 为 $0.729\pm0.009$，RMA 为 $0.996\pm0.011$；Repl-DER 的 A-Dice 为 $0.717\pm0.005$，BWTR 为 $-0.157\pm0.022$。两者均使用较低原始样本回放率。Repl-GSS 的 A-Dice 同为 $0.717\pm0.005$，但 BWTR 为 $-0.276\pm0.055$。这说明即使最终平均性能接近，历史任务退化过程也可能不同。

正则化方法的 A-Dice 位于 0.594 至 0.621 之间，BWTR 均为负。Regu-EWC 和 Regu-SI 的 RMA 接近 1，表明它们能够学习当前器官，但旧任务性能在后续训练中下降。Non-CL 也呈现相似现象，其 RMA 为 $1.006\pm0.008$，A-Dice 仅为 $0.603\pm0.008$。对于标签集合与输入分布共同变化的任务序列，当前任务学习和历史保持之间的差异更为明显。

\begin{table}[htbp]
\centering
\caption{Organ-CL 场景下不同持续学习方法的综合结果}
\label{tab:benchmark-organ-results}
\scriptsize
\setlength{\tabcolsep}{4.5pt}
\renewcommand{\arraystretch}{1.12}
\resizebox{0.90\textwidth}{!}{%
\begin{tabular}{lccccc}
\toprule
方法 & A-Dice & BWTR & RMA & MPE & DRR \\
\midrule
\multicolumn{6}{l}{\textit{正则化方法}} \\
Regu-EWC & $.601\pm.008$ & $-.371\pm.014$ & $1.004\pm.012$ & $0$ & $0$ \\
Regu-SI & $.621\pm.024$ & $-.343\pm.025$ & $1.003\pm.009$ & $0$ & $0$ \\
Regu-LwF & $.594\pm.018$ & $-.308\pm.044$ & $.916\pm.004$ & $0$ & $0$ \\
\addlinespace[2pt]
\multicolumn{6}{l}{\textit{回放方法}} \\
Repl-ER & $.729\pm.009$ & $-.178\pm.021$ & $.996\pm.011$ & $0$ & $.012^{*}$ \\
Repl-GSS & $.717\pm.005$ & $-.276\pm.055$ & $.995\pm.009$ & $0$ & $.023^{*}$ \\
Repl-GEM & $.599\pm.007$ & $-.182\pm.068$ & $.806\pm.059$ & $0$ & $.023^{*}$ \\
Repl-GPM & $.605\pm.009$ & $-.329\pm.034$ & $.960\pm.027$ & $0$ & $.068$ \\
Repl-AGEM & $.635\pm.013$ & $-.193\pm.069$ & $.873\pm.052$ & $0$ & $.023^{*}$ \\
Repl-DER & $.717\pm.005$ & $-.157\pm.022$ & $.953\pm.013$ & $0$ & $.012^{*}$ \\
Repl-FDR & $.631\pm.008$ & $-.217\pm.008$ & $.885\pm.013$ & $0$ & $.023^{*}$ \\
\addlinespace[2pt]
\multicolumn{6}{l}{\textit{参数隔离方法}} \\
ParIso-PNN & $.839\pm.004$ & $0\pm0$ & $1.001\pm.006$ & $1.074$ & $0$ \\
ParIso-DAN & $.824\pm.005$ & $0\pm0$ & $.972\pm.016$ & $.166$ & $0$ \\
\addlinespace[2pt]
\multicolumn{6}{l}{\textit{参照方法}} \\
Non-CL & $.603\pm.008$ & $-.369\pm.005$ & $1.006\pm.008$ & $0$ & $0$ \\
JointTrain & $.881\pm.010$ & -- & -- & $0$ & $0$ \\
\bottomrule
\end{tabular}%
}
\vspace{2pt}
\begin{minipage}{0.90\textwidth}
\footnotesize 注：数值为均值$\pm$标准差；“--”表示不适用或未报告；DRR 中带“$^{*}$”的数值表示回放原始图像。JointTrain 可访问全部任务数据，其结果仅作为不同访问条件下的参照。
\end{minipage}
\end{table}

\subsection*{任务顺序与回放容量分析}

持续学习结果还可能受到任务排列影响。图\ref{fig:benchmark-order-adice}和图\ref{fig:benchmark-order-bwtr}展示了十种 Domain-CL 任务顺序下的 A-Dice 与 BWTR。不同顺序会改变各方法的最终表现，说明固定单一任务序列不足以完全描述方法鲁棒性。Repl-GSS 在十种顺序下的 A-Dice 和 BWTR 标准差分别为 0.0134 和 0.0123，整体波动较小。由于图中未列出全部排列的逐项数值，本章仅比较总体趋势，不进行逐顺序排名。

\begin{figure}[htbp]
\centering
\includegraphics[width=0.92\textwidth]{ch03/task_robustness_adice.pdf}
\caption{十种随机 Domain-CL 任务顺序下各方法的 A-Dice。图中汇总不同任务排列的结果，并以标准差反映方法对任务顺序的敏感性。}
\label{fig:benchmark-order-adice}
\end{figure}

\begin{figure}[htbp]
\centering
\includegraphics[width=0.92\textwidth]{ch03/task_robustness_bwt_from_bottom.pdf}
\caption{十种随机 Domain-CL 任务顺序下各方法的 BWTR。图中标准差反映历史知识保持对任务排列变化的敏感性。}
\label{fig:benchmark-order-bwtr}
\end{figure}

回放容量决定历史样本对训练的覆盖程度，也改变存储成本。图\ref{fig:benchmark-memory}给出了多种回放方法随缓存规模变化的 Domain-CL A-Dice。ER、FDR、GEM、DER 和 GSS 的性能整体随缓存增加而提高，随后呈现趋于饱和的趋势。这一结果说明更大的缓存可以改善历史分布近似，但不能据此确定跨数据集通用的最优缓存容量。缓存收益还需与 DRR、数据治理条件和方法自身的历史信息形式共同解释。

\begin{figure}[htbp]
\centering
\includegraphics[width=0.90\textwidth]{ch03/domain_memory_size_adice.pdf}
\caption{Domain-CL 中不同回放方法的缓存规模与 A-Dice 关系。性能随缓存增加总体提高并逐渐趋于饱和，曲线未显示跨数据集一致的最优容量。}
\label{fig:benchmark-memory}
\end{figure}

\subsection*{跨场景综合分析}

三类场景的结果没有形成脱离资源条件的统一最优方法。参数隔离在多个场景中获得零 BWTR 和较高 A-Dice，但需要扩展模型参数。回放方法通常优于单纯正则化，但其性能依赖缓存大小、回放对象和任务顺序。面向 Class-CL 背景语义变化的 MiB 和 PLOP 能够改善全类别分割与遗忘，却不适用于输出任务完全不同的 Organ-CL。正则化方法保持固定参数量和零回放率，但在域或器官变化较大时容易出现较强遗忘。

结果还说明稳定性、可塑性和前向泛化并不总是一致。EWC、SI 和 Non-CL 在部分场景中具有接近或超过 1 的 RMA，却伴随较负的 BWTR；参数隔离方法可以实现零遗忘，但未来域 E-FWT 不一定同步提高；Class-CL 的 WCD 与 A-Dice 也可能给出不同排序。因此，持续医学图像分割的评价应保留完整阶段—任务性能关系，并把资源条件作为结果解释的一部分。


\section{扩展分析与相关学习范式}
\label{sec:benchmark-extensions}

\subsection{基础分割模型的顺序适配分析}
\label{subsec:benchmark-foundation-model}

前述实验表明，现有持续学习方法在 Domain-CL 中的前向泛化仍然有限。由此产生一个直接问题：大规模预训练获得的通用表征，能否在中心域连续到达时同时改善适应、保持和未来域泛化。为对此进行初步观察，本章进一步在六中心 Domain-CL 上考察 Segment Anything Model（SAM）\cite{kirillov2023sam} 的两种顺序适配方式。需要说明的是，本小节中的 SAM 指基础分割模型 Segment Anything Model，与第五章持续配准方法中使用的锐度感知最小化并非同一概念。

第一种方式从相同的 SAM 初始化出发，按照任务 1 至任务 6 的顺序更新编码器和解码器；第二种方式冻结编码器主体，仅顺序更新低秩适配模块和解码器。低秩适配通过在预训练网络中引入可训练的低秩增量矩阵，减少需要改变的参数范围\cite{hu2022lora,ding2023peft}。该分析不把基础模型视为新的主表基线，而是用于观察强预训练表征和参数高效适配是否会改变阶段--任务性能矩阵的结构。

图\ref{fig:benchmark-sam-lora}给出了两种适配方式的任务性能矩阵。每一行表示完成对应任务训练后的模型，每一列表示测试域；主对角线反映当前域适应，左下三角反映历史域保持，右上三角反映对未来域的直接前向泛化。完整微调在多个主对角线位置获得较高 Dice，说明模型能够适应当前中心，但非对角条目波动较大，且完成后续任务后部分历史域性能明显下降。换言之，基础模型的预训练表示并不会自动消除顺序更新造成的覆盖。

采用 LoRA 后，矩阵整体均值由 0.747 提高到 0.823，非对角条目的平均 Dice 由 0.719 提高到 0.808。提升不仅出现在当前任务，也体现在若干历史域和未来域条目，说明冻结大部分预训练编码器并限制更新范围，有助于保留较稳定的共享表示。不过，该结果不意味着参数高效微调已经解决持续学习问题：解码器和适配参数仍按顺序更新，任务继续增加时仍可能发生遗忘或适配模块累积。更合理的方向是把参数高效适配与正则化、历史子空间保护、适配器选择或合并机制结合，在控制参数增长的同时保留跨域表征。

\begin{figure}[htbp]
\centering
\includegraphics[width=\textwidth]{ch03/SAM_confusion_matrix.pdf}
\caption[基础分割模型在 Domain-CL 中的阶段--任务 Dice 矩阵]{基础分割模型在 Domain-CL 中的阶段--任务 Dice 矩阵。左图为编码器和解码器均参与更新的顺序 SAM 微调，右图为冻结编码器主体、仅更新 LoRA 模块和解码器的顺序适配。行表示完成相应任务后的模型，列表示测试任务；对角线、左下三角和右上三角分别对应当前任务适应、历史任务保持和未来域前向泛化。}
\label{fig:benchmark-sam-lora}
\end{figure}

\subsection{持续学习与相关学习范式的联系}
\label{subsec:benchmark-related-paradigms}

持续医学图像分割中的任务演化、数据访问和标注条件往往同时发生，因此持续学习还会与多种机器学习范式形成交叉。图\ref{fig:benchmark-related-paradigms}概括了这些联系。本章的实证研究仍以监督式、任务边界已知的持续分割为主，图中其他方向用于说明基准结论可以延伸到哪些更复杂的部署问题，而不表示本章已经对所有交叉场景完成实验验证。

\textbf{联邦持续学习。} 联邦学习解决数据在不同中心间不能直接集中共享的问题，持续学习则处理同一中心内部数据或任务随时间变化的问题。二者结合后，模型需要同时面对空间异质性和时间演化，并在不交换原始数据的条件下保留跨客户端历史知识\cite{yang2024fclsurvey,yoon2021fedweit}。传统同步联邦训练还可能受到不同中心任务进度和计算条件不一致的影响。第三章对场景、稳定性、可塑性和资源条件的划分，为第四章进一步研究无回放联邦持续学习提供了评价基础。

\textbf{持续测试时适应。} 测试时适应利用无标签目标数据在部署阶段更新模型，不要求重新访问源训练集\cite{wang2021tent,hu2021fullytta}。这一条件与临床中“新中心数据已经到达，但像素级标注尚未获得”的情况相符。若目标分布继续变化，模型还需要在适应当前测试流的同时保持此前测试域和源域能力，从而形成持续测试时适应问题。本章发现多数方法的前向泛化提升有限，说明仅依赖训练阶段的历史保护不足以保证模型在未知中心直接有效；前向泛化评价、分布变化检测和无标签在线适应需要进一步结合。

图中还列出了少样本持续学习、无监督持续学习、开放世界学习和持续新类发现等交叉方向。这些方向分别放宽当前任务标注充足、任务类别预先给定或任务边界明确等假设。它们可以沿用本章“先明确任务变化和访问条件，再分别评价保持、学习与泛化”的原则，但需要重新定义适合自身输出语义的评价矩阵和指标。

\begin{figure}[htbp]
\centering
\includegraphics[width=0.78\textwidth]{ch03/plots_for_benchmark.pdf}
\caption[持续学习与相关机器学习范式的联系]{持续学习与相关机器学习范式的联系。联邦持续学习、持续测试时适应、少样本持续学习、无监督持续学习和持续新类发现等方向，分别改变数据分布方式、标注可用性或类别开放程度。图示用于说明潜在研究交叉，不表示所有方向均已纳入本章实验。}
\label{fig:benchmark-related-paradigms}
\end{figure}

\section{基准研究的主要发现与局限性}
\label{sec:benchmark-findings-limitations}

\subsection*{主要发现}

第一，医学图像分割中的持续变化不能由单一场景统一描述。跨中心域变化保持标签语义不变，类别增量扩展统一输出空间，跨器官任务则一并改变输入分布和分割目标。不同场景决定哪些性能矩阵条目具有可比性，也决定 E-FWT、WCD 和多头评价是否适用。场景定义必须与标签设定、输出组织和测试方式共同给出。

第二，最终平均性能和遗忘量不足以完整评价持续模型。Domain-CL 和 Organ-CL 中，一些方法的 RMA 接近或超过 1，但 BWTR 仍明显为负，说明当前任务建模与历史保持可以沿不同方向变化。Class-CL 中，WCD 与 A-Dice 对专项方法给出不同排序，表明最终统一输出和逐任务平均需要分别观察。

第三，稳定性改进通常伴随资源或可塑性代价。参数隔离方法能够获得零 BWTR，但 MPE 显示模型容量随任务增长。回放方法利用历史样本改善平均性能和遗忘，却引入非零 DRR，并受缓存规模影响。正则化方法不保存原始数据和扩展参数，但过强约束可能降低当前任务学习。由此，方法比较不能脱离历史访问、参数增长和回放预算。

第四，前向泛化需要独立评价。Domain-CL 中，多数方法的 E-FWT 没有相对于 Non-CL 形成稳定提升；DAN 在表中取得较高均值，但方差较大。降低遗忘不自动意味着模型已经学到对未来中心更有效的表示。E-FWT 通过汇总多个未来域的直接表现补充这一维度，但其定义依赖共享输出语义，不适合直接套用于 Class-CL 和 Organ-CL。

第五，任务顺序和缓存容量会改变方法行为。不同中心排列导致 A-Dice 和 BWTR 波动，缓存增大则总体改善回放方法的结果并逐渐趋于饱和。这些观察说明单一顺序和单一缓存预算只能刻画方法的一部分行为。报告持续学习结果时，应一并说明任务排列、回放容量和输出组织方式。

第六，基础模型提供了更强的预训练起点，但不能替代持续学习机制。Domain-CL 的补充分析中，LoRA 顺序适配相较完整微调改善了总体和非对角 Dice，说明限制参数更新范围有助于保留共享表示；与此同时，历史域退化和适配参数累积仍然存在。基础模型持续适配仍需联合考虑参数效率、历史保持和未来域泛化。

\subsection*{研究局限性}

本章的结论受数据、网络和实验范围限制。三类场景覆盖中心域变化、类别扩展和跨器官任务，但没有研究任务边界未知、连续漂移、中心重访、开放世界类别或多种变化因素叠加发生的情况。主基准中的方法采用 ResUNet32 和固定基础训练预算，结果不能直接代表其他网络或三维模型上的排序。SAM 与 LoRA 的补充分析仅覆盖 Domain-CL 和两种顺序适配方式，也不能据此概括全部医学基础模型的持续学习行为。

实验设置未完整报告各数据集和中心的样本数、随机种子、重复次数、硬件、早停规则及代码版本。结果表提供均值与标准差，但缺少相应统计检验，因此本章主要比较数值趋势。Class-CL 的旧类标签访问、背景处理和任务身份设置也限制了与其他标签设定的直接对照。

参数和数据资源评价仍不完整。MPE 反映参数增长，DRR 反映回放比例，但没有覆盖训练时间、推理延迟、显存、能耗和通信成本。JointTrain 使用全部历史数据，只能作为不同访问条件下的参照。未来工作还需要在统一代码、多个任务顺序和更完整资源核算下复核本章观察，并扩展至 任务身份未知（task-free）、开放世界和基础模型持续适配场景。


\section{本章小结}
\label{sec:benchmark-summary}

本章围绕持续医学图像分割的场景定义与综合评价构建了系统性基准。根据输入分布、标签空间和分割对象的变化，基准组织 Domain-CL、Class-CL 和 Organ-CL 三类场景，并在统一主干、训练预算和数据划分下比较正则化、回放、参数隔离及类别增量专项方法。三类场景服务于本章的数据组织和实验分析，不被视为对全部持续医学影像问题的穷尽分类。

在评价方面，本章在 A-Dice 和 BWTR 的基础上，引入 RMA 描述持续约束下的当前任务建模能力，使用 E-FWT 评价 Domain-CL 对多个未来域的直接表现，并以 WCD、MPE 和 DRR 补充分类别输出及资源条件。实验结果表明，知识保持、当前任务学习、前向泛化和资源开销可能给出不同方法排序；任务顺序和回放容量也会改变方法行为。基础分割模型的补充实验进一步表明，预训练表征和参数高效适配能够改善顺序域学习的性能矩阵，但不会自动消除遗忘。联邦持续学习和持续测试时适应等交叉方向则要求进一步放宽数据集中、标注及时可用等假设。由此，后续持续医学影像研究需要在明确任务设定和访问条件的基础上解释结果，而不能只依据单一最终性能或遗忘指标。第三章建立的场景与评价语言将作为第四章联邦持续分类和第五章任务级持续学习研究的参照，但不构成后续方法的算法组件。
